% Part: sets-functions-relations
% Chapter: infinite
% Section: dedekind
%
\documentclass[../../../include/open-logic-section]{subfiles}

\begin{document}

\olfileid{sfr}{infinite}{dedekind}
\olsection{ડેડેકિન્ડ બીજગણિતો}

આપણે પ્રાકૃતિક સંખ્યાઓ માત્ર અનંત હોય એટલું જ નથી માગતા; આપણે એમ પણ
માગીએ છીએ કે તેમને અમુક (બૈજિક) ગુણધર્મો હોય: સરવાળા, ગુણાકાર વગેરે
હેઠળ તેમનો વ્યવહાર સુસંગત હોવો જોઈએ.

ડેડેકિન્ડનો વિચાર \emph{અનુગામી વિધેય}ના ખ્યાલને મૂળભૂત માનવાનો અને
પછી નીચેના ગુણધર્મો ધરાવતી સંખ્યાઓ તરીકે તેમનું લક્ષણ આપવાનો હતો:
\begin{enumerate}
	\item એવી સંખ્યા $0$ છે જે કોઈ પણ સંખ્યાની અનુગામી નથી
	\\અર્થાત્, $0 \notin \ran{s}$
	\\અર્થાત્, $\forall x\ s(x) \neq 0$
	\item ભિન્ન સંખ્યાઓના અનુગામીઓ ભિન્ન હોય છે
	\\અર્થાત્, $s$ એ !!a{injection} છે
	\\અર્થાત્, $\forall x \forall y (s(x) = s(y) \lif x = y)$
	\item\ollabel{repeatedapplication} દરેક સંખ્યા $0$ પર અનુગામી
	વિધેય વારંવાર લગાડવાથી મળે છે.
\end{enumerate}
પહેલી બે શરતોને પ્રથમ-ક્રમ તર્કશાસ્ત્રથી સહેલાઈથી સંભાળી શકાય છે
(ઉપર જુઓ). પરંતુ માત્ર પ્રથમ-ક્રમ તર્કશાસ્ત્રથી
\olref{repeatedapplication}ને સંભાળી શકાતી નથી. ડેડેકિન્ડની સફળતા
શરત \olref{repeatedapplication}ને ગણસૈદ્ધાંતિક રીતે આ પ્રમાણે ફરી
ઘડવામાં હતી:
\begin{enumerate}
	\item[3$'$.] પ્રાકૃતિક સંખ્યાઓ એવો નાનામાં નાનો ગણ છે જે
	\emph{અનુગામી વિધેય હેઠળ સંવૃત} છે: એટલે કે, ગણના કોઈ પણ
	!!{element} પર $s$ લગાડીએ તો ગણનો બીજો !!{element} મળે છે.
\end{enumerate}
પરંતુ આ વાત આપણે ધીમે ધીમે પૂરેપૂરી સ્પષ્ટ કરવી પડશે.

\begin{defn}\ollabel{Closure}
	દરેક વિધેય $f \colon A \to A$ માટે ગણ $X \subseteq A$ એ
	$f$-\emph{સંવૃત} છે {iff} $(\forall x \in X)f(x) \in X$.
	હવે દરેક $o \in A$ માટે વ્યાખ્યા કરીએ:
	$$\closureofunder{f}{o} = \bigcap\Setabs{X}{o \in X\text{ અને }X
\text{ એ $f$-સંવૃત છે}}$$
\end{defn}
\sourcecorrection{OLINF-001}{મૂળની \texttt{dedekind-algebra.tex},
પંક્તિઓ 40--64માં સંવરણની વ્યાખ્યા કોઈ પણ $f$ અને $o$ માટે અપાઈ
છે, પછીના લેમામાં $A$ને બાંધ્યા વિના $o\in A$ લખાયું છે, અને
સાબિતીમાં $\ran{f}\cup\{o\}$ એ $f$-સંવૃત હોવાનો ઉપયોગ થાય છે.
આ દલીલને જરૂરી અને પછીના બધા ઉપયોગો સાથે સુસંગત એવી ધારણા
$f\colon A\to A$, $X\subseteq A$ અને $o\in A$ અહીં સ્પષ્ટ કરી છે.}

આમ $\closureofunder{f}{o}$ એ એવા બધા $f$-સંવૃત ગણોનો છેદ છે
જેમનો !!a{element}~$o$ છે. તેથી સહજ રીતે,
$\closureofunder{f}{o}$ એ !!a{element} તરીકે $o$ ધરાવતો
\emph{નાનામાં નાનો} $f$-સંવૃત ગણ છે. આગળનું પરિણામ આ સહજ વિચારને
ચોક્કસ બનાવે છે;
\begin{lem}\ollabel{closureproperties}
	દરેક વિધેય $f \colon A \to A$ અને દરેક $o \in A$ માટે:
	\begin{enumerate}
		\item\ollabel{closurehaselem} $o \in \closureofunder{f}{o}$; અને
		\item\ollabel{closureclosed} $\closureofunder{f}{o}$ એ $f$-સંવૃત છે; અને
		\item\ollabel{closuresmallest} જો $X$ એ $f$-સંવૃત હોય અને $o \in
		X$, તો $\closureofunder{f}{o} \subseteq X$.
	\end{enumerate}
\end{lem}

\begin{proof}
નોંધો કે !!a{element} તરીકે $o$ ધરાવતો ઓછામાં ઓછો એક $f$-સંવૃત ગણ
છે, એટલે કે $\ran{f}\cup \{o\}$. તેથી આવા \emph{બધા} ગણોનો છેદ
$\closureofunder{f}{o}$ અસ્તિત્વ ધરાવે છે. હવે આપણે
\olref{closurehaselem}--\olref{closuresmallest} તપાસવાં છે.

\olref{closurehaselem} અંગે: $o \in \closureofunder{f}{o}$, કારણ કે
તે એવા ગણોનો છેદ છે જેમના દરેકનો !!a{element}~$o$ છે.

\olref{closureclosed} અંગે: ધારો કે $x \in \closureofunder{f}{o}$.
તેથી જો $o \in X$ અને $X$ એ $f$-સંવૃત હોય, તો $x \in X$; અને
$X$ એ $f$-સંવૃત હોવાથી હવે $f(x) \in X$. તેથી
$f(x) \in \closureofunder{f}{o}$.

\olref{closuresmallest} અંગે: સામાન્ય રીતે, જો $X \in C$ હોય, તો
$\bigcap C \subseteq X$.
\end{proof}

આનો ઉપયોગ કરીને આપણે કહી શકીએ:

\begin{defn}
\emph{ડેડેકિન્ડ બીજગણિત} એટલે ગણ $A$, વિધેય $f \colon A \to A$
અને કોઈ $o \in A$ એવાં કે:
	\begin{enumerate}
		\item \ollabel{ded:proper} $o \notin \ran{f}$
		\item \ollabel{ded:injection} $f$ એ !!a{injection} છે
		\item \ollabel{ded:closure} $A = \closureofunder{f}{o}$
	\end{enumerate}
\end{defn}

$A = \closureofunder{f}{o}$ હોવાથી આપણું પહેલાનું પરિણામ જણાવે છે કે
$A$ એ !!a{element} તરીકે $o$ ધરાવતો નાનામાં નાનો $f$-સંવૃત ગણ છે.
ડેડેકિન્ડ બીજગણિત સ્પષ્ટપણે ડેડેકિન્ડ-અનંત છે; વ્યાખ્યાની કલમો
\olref{ded:proper} અને \olref{ded:injection} જ જુઓ. પરંતુ વધુ રસપ્રદ
હકીકત એ છે કે કોઈ પણ ડેડેકિન્ડ-અનંત ગણને ડેડેકિન્ડ બીજગણિતમાં ફેરવી
શકાય છે.

\begin{thm}\ollabel{thm:DedekindInfiniteAlgebra}
જો કોઈ ડેડેકિન્ડ-અનંત ગણ હોય, તો કોઈ ડેડેકિન્ડ બીજગણિત હોય છે.
\end{thm}

\begin{proof}
ધારો કે $D$ ડેડેકિન્ડ-અનંત છે. તેથી એક-એક વિધેય $g \colon D \to
D$ અને ઘટક $o \in D \setminus \ran{g}$ છે. હવે $A =
\closureofunder{g}{o}$ લો; \olref{closureproperties} મુજબ $A$
અસ્તિત્વ ધરાવે છે અને $o \in A$. $f = \funrestrictionto{g}{A}$ લો.
$A, f, o$ મળીને ડેડેકિન્ડ બીજગણિત રચે છે તે આપણે બતાવીશું.

\olref{ded:proper} અંગે: $o \notin \ran{g}$ અને $\ran{f}
\subseteq \ran{g}$ હોવાથી $o\notin \ran{f}$.

\olref{ded:injection} અંગે: $g$ એ $D$ પર એક-એક વિધેય છે; તેથી
$f \subseteq g$ પણ એક-એક વિધેય જ હોય.

\olref{ded:closure} અંગે: \olref{closureproperties} મુજબ $A$ એ
$g$-સંવૃત છે; તેથી ખાસ કરીને $A$ એ $f$-સંવૃત છે. આમ
$\closureofunder{f}{o} \subseteq A$ એ \olref{closureproperties}થી મળે
છે. વળી $\closureofunder{f}{o}$ એ $f$-સંવૃત છે અને
$f = \funrestrictionto{g}{A}$ હોવાથી $\closureofunder{f}{o}$ એ
$g$-સંવૃત છે. તેથી \olref{closureproperties} મુજબ
$A \subseteq \closureofunder{f}{o}$.
\end{proof}

\end{document}
